% Generated draft: requires mathematical and source review.
\documentclass[11pt]{article}
\usepackage[margin=25mm]{geometry}
\usepackage{fontspec}
\setmainfont{Latin Modern Roman}
\usepackage{amsmath,amssymb,mathtools,unicode-math}
\setmathfont{Latin Modern Math}
\usepackage{enumitem,needspace,xurl,hyperref}
\setlength{\parindent}{0pt}
\setlength{\parskip}{5pt}
\setlength{\emergencystretch}{3em}
\newcommand{\N}{\mathbb{N}}
\newcommand{\Z}{\mathbb{Z}}
\newcommand{\Q}{\mathbb{Q}}
\newcommand{\R}{\mathbb{R}}
\newcommand{\C}{\mathbb{C}}
\newcommand{\F}{\mathbb{F}}
\newcommand{\A}{\mathbb{A}}
\newcommand{\PP}{\mathbb{P}}
\newcommand{\E}{\mathbb{E}}
\newcommand{\eps}{\varepsilon}
\newcommand{\dd}{\,\mathrm{d}}
\begin{document}
\section*{UnsolvedMath ID: 701111}
\textbf{Source problem code:} catalogue entry
\subsection{Definitions and mathematical statement}
The supplied record does not define \emph{barely transitive}. The formulation below uses the convention that this means a faithful transitive action on an infinite set in which every proper subgroup has only finite orbits. This interpretation, including the infinitude requirement, requires confirmation against the cited source.

Let \(\Omega\) be an infinite set and let \(G\leq\operatorname{Sym}(\Omega)\). The action is \emph{transitive} if for all \(\alpha,\beta\in\Omega\) there exists \(g\in G\) with \(g\alpha=\beta\). A \emph{block} is a nonempty subset \(B\subseteq\Omega\) such that, for every \(g\in G\), either \(gB=B\) or \(gB\cap B=\varnothing\). As defined in the supplied statement, the action is \emph{totally imprimitive} if
\[
\forall S\subseteq\Omega\text{ finite},\quad
\exists B\subseteq\Omega\text{ a finite block with }S\subseteq B.
\]
Under the convention stated above, the action is \emph{barely transitive} if it is transitive and
\[
\forall H<G\;\forall\alpha\in\Omega,\qquad
|H\alpha|<\infty,
\qquad H\alpha=\{h\alpha:h\in H\}.
\]
A group is \emph{locally finite} if every subgroup generated by a finite subset is finite.

\textbf{Question (V. V. Belyaev, 2002; Kourovka Notebook Problem 15.23).} Does there exist an infinite set \(\Omega\) and a transitive group \(G\leq\operatorname{Sym}(\Omega)\) whose action is totally imprimitive and barely transitive, but for which
\[
\exists F\subseteq G\text{ finite},\qquad |\langle F\rangle|=\infty?
\]
The supplied registry labels this problem open; that label is not an independently verified assessment.
\subsection{Short English statement}
Under the stated interpretation of barely transitive, the question asks whether an infinite transitive permutation group can have every finite set of points contained in a finite block and every proper subgroup acting with finite orbits, while possessing an infinite finitely generated subgroup. The interpretation of barely transitive requires source confirmation.
\subsection{Sources}
\begin{itemize}
\item Local UnsolvedMath registry, permanent ID \(701111\), exact source code \texttt{catalogue entry}, title \emph{Kourovka Notebook Problem 15.23}; supplied upstream revision \texttt{d371ee917ac9647fbb06cd6fc529deec6f3dda82}. The supplied status is \texttt{open}.
\item \emph{The Kourovka Notebook (Unsolved Problems in Group Theory)}, 15th issue (2002), Problem 15.23, proposed by V. V. Belyaev, as attributed in the supplied record.
\item The supplied background also identifies the 21st issue edition (2026), \texttt{arXiv:1401.0300v48}. The collection description credits E. I. Khukhro and V. D. Mazurov as editors of Issues 1--21 (1965--2026). These bibliographic details are supplied attribution and have not been independently checked.
\end{itemize}
\end{document}
